\documentclass[11pt,a4paper]{article}

\usepackage[utf8]{inputenc}
\usepackage[T1]{fontenc}
\usepackage[english]{babel}
\usepackage[a4paper,margin=2.3cm]{geometry}
\usepackage{amsmath,amssymb,amsfonts,amsthm}
\usepackage{mathtools}
\usepackage{bm}
\usepackage{booktabs}
\usepackage{enumitem}
\usepackage{xcolor}
\usepackage{hyperref}
\usepackage{tcolorbox}

\hypersetup{
  colorlinks=true,
  linkcolor=blue!50!black,
  urlcolor=blue!60!black,
  citecolor=blue!40!black
}

\newcommand{\BHF}{B_{\mathrm{hf}}}
\newcommand{\dEQ}{\Delta E_{Q}}
\newcommand{\code}[1]{\texttt{\small #1}}
\newcommand{\file}[1]{\textcolor{black!70}{\texttt{\small #1}}}
\DeclareMathOperator{\diag}{diag}
\DeclareMathOperator{\tr}{tr}
\DeclareMathOperator{\erfc}{erfc}
\DeclareMathOperator*{\argmin}{arg\,min}
\DeclareMathOperator*{\argmax}{arg\,max}

\newtcolorbox{nota}[1][Note]{
  colback=blue!4, colframe=blue!45!black, fonttitle=\bfseries,
  title={#1}, left=4pt, right=4pt, top=3pt, bottom=3pt
}

\title{\textbf{Mathematical manual of the \emph{Mossbauer} fitting engine}\\[3pt]
\large Forward model, discrete fitting and distribution fitting}
\author{Department of Physical Chemistry · UAM}
\date{\today}

\begin{document}
\maketitle

\begin{abstract}\noindent
This manual documents, in a self-contained way, the complete mathematics of the
fitting engine for $^{57}$Fe Mössbauer spectra implemented in the application. It
covers (i)~the \emph{forward model} —line profiles, positions and intensities of the
six lines of the sextet with first-order and Kündig treatments of the
magnetic–quadrupole coupling, doublets and singlets—; (ii)~the \emph{discrete
fitting} —statistics of the data in Gauss and Poisson modes, propagation of the
calibration uncertainty, cost functional with robust losses,
trust-region optimization with multistart and a global option,
covariance, diagnostics and bootstrap—; and (iii)~the \emph{distribution fitting}
$P(\BHF)$ or $P(\dEQ)$ —discretization, kernel, Tikhonov and total
variation regularization, conditioning, effective degrees of freedom, automatic
selection of the regularization parameter, error bars and parametric
distributions. References to the file and function of the code are included.
\end{abstract}

\tableofcontents
\bigskip

%======================================================================
\part{Forward model}
%======================================================================

\section{Transmission spectrum}\label{sec:fwd}

After folding the raw data about the folding point, the spectrum is
defined over $N$ channels with velocities $\{v_i\}_{i=1}^{N}$ equally spaced in
$[-v_{\max}, v_{\max}]$ and normalized transmission $\{y_i\}$, with $y_i\approx 1$
away from the absorption peaks. The forward model is a linear superposition of
absorptions over a baseline:
\begin{equation}
  \boxed{\;
  T(v;\bm\theta) \;=\; b + s\,v \;-\; \sum_{c=1}^{N_c} A_c(v;\bm\theta_c)
  \;}
  \label{eq:Tmodel}
\end{equation}
where $b$ is the baseline (multiplicative, the data are normalized by the
\code{norm\_factor}), $s$ the slope, and each $A_c$ is the absorption of a
component (sextet, doublet or singlet). It is a \emph{thin absorber}
model: linear in the depth of each component. Implemented in
\file{core/physics.py}, function \code{total\_model}.

\section{Line profiles}\label{sec:lineshape}

Each elementary line uses a profile normalized to unit height at its center. The
application offers two profiles, selectable globally.

\subsection{Lorentzian}
\begin{equation}
  L(v;v_0,\gamma) \;=\; \frac{\gamma^{2}}{(v-v_0)^{2}+\gamma^{2}},
  \label{eq:lorentz}
\end{equation}
with $\gamma$ the half-width at half maximum (HWHM). It satisfies $L(v_0)=1$.

\subsection{Pseudo-Voigt (Faddeeva)}
The convolution of a Lorentzian (HWHM $\gamma$) with a Gaussian (width
$\sigma$) is the Voigt profile, evaluated through the Faddeeva
function $w(z)=e^{-z^2}\erfc(-iz)$:
\begin{equation}
  \widetilde V(v;v_0,\gamma,\sigma) \;=\;
   \frac{1}{\sigma\sqrt{2\pi}}\,\Re\!\left[w\!\left(\frac{(v-v_0)+i\gamma}{\sigma\sqrt2}\right)\right].
\end{equation}
So that the profile has height exactly $1$ at $v_0$ \emph{independently
of the sampling}, it is normalized by its analytic value at the peak:
\begin{equation}
  V(v;v_0,\gamma,\sigma) \;=\; \frac{\widetilde V(v;v_0,\gamma,\sigma)}{\widetilde V(v_0;v_0,\gamma,\sigma)},
  \qquad
  \widetilde V(v_0) = \frac{1}{\sigma\sqrt{2\pi}}\,\Re\!\left[w\!\left(\tfrac{i\gamma}{\sigma\sqrt2}\right)\right].
  \label{eq:voigt}
\end{equation}
Since $w(iy)=e^{y^2}\erfc(y)\in\mathbb R^{+}$ for $y\in\mathbb R$, the denominator
is real and positive. In the limit $\sigma\to0$, $V\to L$.

\begin{nota}[Normalization to the peak]
Normalizing by the \emph{discrete} maximum of the grid (as previous versions
did) underestimates the real peak when no $v_i$ falls on $v_0$, inflating the
amplitude and making the depth depend on the sampling. The form
\eqref{eq:voigt} removes that bias.
\end{nota}

\section{Magnetic sextet}\label{sec:sextet}

A sextet arises from the nuclear Zeeman splitting combined, when applicable, with the
quadrupole interaction. The six lines correspond to the
allowed transitions ($\Delta m=0,\pm1$) between the ground state $I_g=\tfrac12$ and the
excited state $I_e=\tfrac32$. The absorption is
\begin{equation}
  A_{\text{sext}}(v) \;=\; d\,\sum_{j=1}^{6} W_j\, P\!\left(v;v_{0,j},\gamma_j\right),
  \label{eq:Asext}
\end{equation}
with $d$ the depth, $P$ the profile (\S\ref{sec:lineshape}), $W_j$ the
intensities (\S\ref{sec:intensities}) and $\gamma_j$ the linewidths
\begin{equation}
  \bm\gamma = \gamma_1\,[\,1,\;\gamma_2,\;\gamma_3,\;\gamma_3,\;\gamma_2,\;1\,],
\end{equation}
where $\gamma_2,\gamma_3$ are \emph{relative} linewidths. The positions
$v_{0,j}$ are obtained by one of two treatments.

\subsection{First-order treatment}
The quadrupole is treated as a rigid perturbation on the Zeeman pattern:
\begin{equation}
  v_{0,j} \;=\; \frac{\BHF}{B_0}\,r_j^{(0)} \;+\; \delta \;+\; s_j\,\dEQ,
  \qquad j=1,\dots,6,
  \label{eq:sext_pos1}
\end{equation}
with $B_0=33\,$T the reference field, $\delta$ the isomer shift,
$r^{(0)}$ the positions calibrated at $B_0$ (Appendix~\ref{ap:calib}) and the
quadrupole pattern
\begin{equation}
  \bm s = \left[\,+\tfrac12,\,-\tfrac12,\,-\tfrac12,\,-\tfrac12,\,-\tfrac12,\,+\tfrac12\,\right].
\end{equation}
Valid when $\dEQ \ll \mu_N g_e \BHF$ and the axis of the electric field gradient
(EFG) is parallel to $\vec B$.

\subsection{Kündig treatment (full axial Hamiltonian)}\label{sec:kundig}
For an axial EFG (asymmetry parameter $\eta=0$) and an arbitrary angle $\beta$ between
$\vec B$ and the principal axis $V_{zz}$, the positions are obtained by diagonalizing the
Hamiltonian of the excited state. In the basis $\{\lvert\tfrac32\rangle,
\lvert\tfrac12\rangle, \lvert{-}\tfrac12\rangle, \lvert{-}\tfrac32\rangle\}$:
\begin{equation}
  \mathcal H_e \;=\; \omega_e\, I_z \;+\; \frac{\dEQ}{6}\bigl(3 I_{z'}^2 - I(I+1)\bigr),
  \qquad I_{z'} = I_z\cos\beta + I_x\sin\beta,
  \label{eq:Hkundig}
\end{equation}
with $I(I+1)=\tfrac{15}{4}$ and $\omega_e=(\BHF/B_0)\,\omega_e^{(0)}$ the
Zeeman splitting of the excited state. The numerical diagonalization of the
$4\times4$ matrix gives the eigenvalues $E^{e}_{m}$ and eigenvectors. The ground state
($I_g=\tfrac12$) is purely Zeeman:
\begin{equation}
  E^{g}_{\pm 1/2} = \pm\tfrac12\,\omega_g, \qquad \omega_g=(\BHF/B_0)\,\omega_g^{(0)}.
\end{equation}
The six positions are the energy differences of the allowed transitions,
assigning to each excited-state eigenvalue its dominant $m_e$ (by the overlap of
the eigenvector):
\begin{equation}
\begin{aligned}
  v_{0,1}&=E^{e}_{-3/2}-E^{g}_{-1/2}, & v_{0,2}&=E^{e}_{-1/2}-E^{g}_{-1/2}, & v_{0,3}&=E^{e}_{+1/2}-E^{g}_{-1/2},\\
  v_{0,4}&=E^{e}_{-1/2}-E^{g}_{+1/2}, & v_{0,5}&=E^{e}_{+1/2}-E^{g}_{+1/2}, & v_{0,6}&=E^{e}_{+3/2}-E^{g}_{+1/2},
\end{aligned}
\end{equation}
plus $\delta$. In the limit $\beta=0$ it reproduces \eqref{eq:sext_pos1} to machine
precision. Implemented in \file{core/hamiltonian.py}.

\subsection{Polycrystalline average}\label{sec:powder}
In a powder sample all orientations $\beta$ are equally probable. The
absorption is averaged over the sphere:
\begin{equation}
  A_{\text{sext}}^{\text{pol}}(v) \;=\; \frac12\int_0^{\pi}\!\sin\beta\;A_{\text{sext}}(v;\beta)\,d\beta
  \;\approx\; \sum_{k=1}^{N_q} w_k\, A_{\text{sext}}(v;\beta_k),
  \label{eq:powder}
\end{equation}
with Gauss–Legendre quadrature nodes and weights over $x=\cos\beta\in[-1,1]$,
$\beta_k=\arccos x_k$, $w_k=\tfrac12 w_k^{\mathrm{GL}}$ (so that
$\sum_k w_k=1$). The computation is vectorized: for $N_q$ orientations a
tensor $(N_q,4,4)$ is built and diagonalized at once. Default value
$N_q=20$ (error $<0.01\%$).

\section{Intensities and texture}\label{sec:intensities}

The intensities of the six lines respect the symmetry $1\!\leftrightarrow\!6$,
$2\!\leftrightarrow\!5$, $3\!\leftrightarrow\!4$:
\begin{equation}
  \bm W = [\,I_1,\,I_2,\,I_3,\,I_3,\,I_2,\,I_1\,],\qquad
  I_1=i_3 i_1,\;\; I_2=i_3 i_2,\;\; I_3=i_3.
\end{equation}
The application offers two parametrizations per sextet.

\subsection{Free intensities}
$i_1,i_2,i_3$ are independent fit parameters. It does not impose the
$3{:}2{:}1$ ratio; useful for numerical flexibility and empirical interpretation.

\subsection{Texture parameter}\label{sec:texture}
For a sample with angle $\theta$ between the magnetic moment and the direction of the
$\gamma$ photon, the magnetic dipole amplitudes are
\begin{equation}
  W_{1,6}\propto 3(1+\cos^2\theta),\quad W_{2,5}\propto 4\sin^2\theta,\quad W_{3,4}\propto (1+\cos^2\theta).
\end{equation}
Defining the texture parameter $t=\sin^2\theta\in[0,1]$ and normalizing by
$(1+\cos^2\theta)=2-t$:
\begin{equation}
  \boxed{\;i_1 = 3,\qquad i_2 = \frac{4t}{2-t},\qquad i_3 = 1.\;}
  \label{eq:texture}
\end{equation}
Cases: $t=\tfrac23\Rightarrow 3{:}2{:}1$ (random powder); $t=1\Rightarrow
3{:}4{:}1$ (planar texture, foil $\perp$ beam); $t=0\Rightarrow 3{:}0{:}1$ (easy
axis $\parallel$ beam). In this mode $t$ is the only free
intensity parameter; $i_1,i_2,i_3$ are derived from \eqref{eq:texture}.

\section{Doublet and singlet}

\begin{align}
  A_{\text{dob}}(v) &= d\bigl[\,i_1\,P(v;\delta-\tfrac{\dEQ}{2},\gamma_1) + i_1 i_2\,P(v;\delta+\tfrac{\dEQ}{2},\gamma_1\gamma_2)\,\bigr],\\
  A_{\text{sing}}(v) &= d\,i_1\,P(v;\delta,\gamma_1).
\end{align}

%======================================================================
\part{Discrete fitting}
%======================================================================

\section{Parametrization and constraints}\label{sec:disc-param}

The vector of free parameters $\bm x\in\mathbb R^{p}$ is composed of the global ones
$\{b,s\}$ (plus optionally $v_{\max}$ and the folding center $c_{\text{fold}}$) and,
for each active component, a subset of
$\bm\theta_c=(\delta,\dEQ,\BHF,\gamma_1,\gamma_2,\gamma_3,d,i_1,i_2,i_3,t,\beta)$
according to its type and mode:
\begin{itemize}[leftmargin=*]
  \item \emph{Type}: singlet uses $\{\delta,\gamma_1,d,i_1\}$; doublet adds
        $\{\dEQ,\gamma_2,i_2\}$; sextet uses all of them.
  \item \emph{Intensities}: in texture mode, $\{i_1,i_2,i_3\}$ are replaced by
        the single parameter $t$.
  \item \emph{Quadrupole}: in Kündig mode with fixed $\beta$, $\beta$ enters as a
        free parameter; in polycrystal or first-order mode, it does not.
  \item \emph{Fixed}: marked parameters are removed from the active vector.
  \item \emph{Linearly constrained} $k_t=\alpha\,k_s+\beta_0$: applied at each
        evaluation (up to six iterations for chains of dependencies).
\end{itemize}
Each parameter has hard bounds $[\ell_k,u_k]$ (Appendix~\ref{ap:bounds}).

\section{Statistics of the data}\label{sec:weighting}

\subsection{Gauss mode}
Folding averages two independent Poisson channels; for counts $c_i$ per
folded channel, $\mathrm{Var}\approx c_i/2$, from which the $1\sigma$ uncertainty
of the normalized datum:
\begin{equation}
  \sigma_i \;=\; \frac{\max\!\bigl(\sqrt{c_i/2},\,1\bigr)}{f_{\text{norm}}},
  \qquad \sigma_i\leftarrow\max(\sigma_i,10^{-9}).
  \label{eq:sigma-gauss}
\end{equation}

\subsection{Poisson mode (IRLS)}
The Poisson likelihood is approximated by reweighted least squares: the
uncertainty is evaluated on the \emph{model} (predicted counts
$\hat c_i = T_i f_{\text{norm}}$) and is recomputed at each evaluation of the
optimizer, which is equivalent to a Newton step on the Poisson $-2\log\mathcal L$:
\begin{equation}
  \sigma_i^{\text{Poiss}} \;=\; \frac{\sqrt{\max(\hat c_i/2,\,1)}}{f_{\text{norm}}}
  \;=\; \frac{\sqrt{\max(T_i f_{\text{norm}}/2,\,1)}}{f_{\text{norm}}}.
  \label{eq:sigma-poisson}
\end{equation}

\subsection{Propagation of the calibration uncertainty}\label{sec:calib}
If the velocity calibration contributes an uncertainty $\sigma_{v_{\max}}$, its
effect on the model is added in quadrature. Since $v_i=v_{\max}(2i/(N-1)-1)$, one
has $\partial v_i/\partial v_{\max}=v_i/v_{\max}$, and by the chain rule
\begin{equation}
  \sigma_i^{\text{eff}\,2} \;=\; \sigma_i^2 \;+\;
  \left(\frac{\partial T}{\partial v}\bigg|_i \frac{v_i}{v_{\max}}\,\sigma_{v_{\max}}\right)^{\!2}.
  \label{eq:calib}
\end{equation}
The term weighs more on the flanks of the peaks (where $\partial T/\partial v$ is
large), which is where a scale error biases $\delta$ and $\BHF$.

\section{Cost functional and robust losses}\label{sec:cost}

The weighted residual is
\begin{equation}
  r_i(\bm x) \;=\; \frac{T(v_i;\bm x)-y_i}{\sigma_i^{\text{eff}}},
  \label{eq:residuo}
\end{equation}
and the minimized functional
\begin{equation}
  \mathcal C(\bm x) \;=\; \tfrac12 \sum_{i=1}^{N}\rho\!\left(r_i(\bm x)^2\right),
  \label{eq:cost}
\end{equation}
with $\rho$ the loss function: \emph{linear} $\rho(z)=z$ (least squares),
\emph{soft-}$\ell_1$ $\rho(z)=2(\sqrt{1+z}-1)$, or \emph{Huber}. The robust ones use
scale $f_{\text{scale}}=3$ (threshold $\sim3\sigma$), reducing the weight of channels
with spurious peaks. For linear $\rho$, $\mathcal C=\tfrac12\chi^2$.

\section{Optimization}\label{sec:opt}

\subsection{Trust region (TRF)}
Minimization is performed with \emph{Trust Region Reflective} over the bounds:
\begin{equation}
  \hat{\bm x} \;=\; \argmin_{\bm\ell\le\bm x\le\bm u}\; \mathcal C(\bm x),
\end{equation}
solving at each iteration the quadratic subproblem with a finite-difference
Jacobian:
\begin{equation}
  \min_{\bm d}\;\tfrac12\|J_k\bm d + \bm r_k\|^{2}
  \;\;\text{s.t.}\;\; \|\bm D_k\bm d\|\le\Delta_k,\;\; \bm\ell-\bm x_k\le\bm d\le\bm u-\bm x_k.
\end{equation}

\subsection{Multistart}
The user's point plus 8 Gaussian restarts are tried:
\begin{equation}
  \bm x^{(s)}=\bm x^{(0)}+\bm\xi^{(s)},\quad \xi^{(s)}_k\sim\mathcal N(0,\rho_k^2(u_k-\ell_k)^2),
\end{equation}
with $\rho_k=0.12$ for $\{\delta,\dEQ,\BHF,\gamma_1,d,v_{\max}\}$ and $0.08$ for the
rest; fixed seed. The one with the lowest cost is retained.

\subsection{Global optimization (optional)}
As an optional pre-pass, \emph{differential evolution} is run over the scalar
cost $\tfrac12\|\bm r\|^2$ (Sobol initialization, no polishing), and its best point
is added as a seed for the TRF polishing. It covers local minima from label
permutation and $\BHF_1\!\leftrightarrow\!\BHF_2$ swaps in
multi-sextet spectra that the Gaussian multistart does not reach.

\section{Uncertainties and diagnostics}\label{sec:disc-stats}

\subsection{Covariance}
From the final Jacobian $J=U\Sigma V^\top$ (SVD, truncated singular values):
\begin{equation}
  \boxed{\;\widehat{\mathrm{Cov}}(\hat{\bm x}) = V\,\Sigma^{-2}V^\top\cdot\frac{2\,\mathcal C(\hat{\bm x})}{N-p},
  \qquad \hat\sigma_{x_k}=\sqrt{\widehat{\mathrm{Cov}}_{kk}}.\;}
  \label{eq:cov}
\end{equation}
Correlations $|\rho_{kl}|\ge0.95$ are reported as a warning of degeneracy.

\subsection{Goodness-of-fit statistics}
\begin{align}
  \chi^2 &= \sum_i r_i^2, & \nu&=N-p, & \chi^2_\nu&=\chi^2/\nu,\\
  \mathrm{AIC}&=N\ln(\mathrm{RSS}/N)+2p, & \mathrm{BIC}&=N\ln(\mathrm{RSS}/N)+p\ln N.
\end{align}

\subsection{Residual diagnostics}
To detect non-random structure (with $\tilde r=r-\bar r$):
\begin{align}
  \text{lag-1 autocorrelation:}\quad &\rho_1=\frac{\sum_i\tilde r_i\tilde r_{i+1}}{\sum_i\tilde r_i^2},\\
  \text{runs (Wald–Wolfowitz):}\quad &Z=\frac{R-\mu_R}{\sigma_R},\\
  \text{antisymmetric correlation:}\quad &c_{\text{as}}=\frac{\tilde{\bm r}\cdot(-\tilde{\bm r}^{\text{rev}})}{\tilde{\bm r}\cdot\tilde{\bm r}},
\end{align}
with warnings if $|\rho_1|>0.35$, $|Z|>2$ or $c_{\text{as}}>0.45$ (the latter
detects centering/calibration errors).

\subsection{Bootstrap}
Monte Carlo estimation of errors: $M$ replicas of synthetic data
$\tilde y^{(m)}=T(\cdot;\hat{\bm x})+\eta^{(m)}$, with $\eta^{(m)}$ Gaussian of
standard deviation $\sigma_i$ (or real Poisson resampling
$\tilde c\sim\mathrm{Poisson}(\hat c)$ in Poisson mode); each replica is refitted and
$\hat\sigma_{x_k}^{\text{boot}}=\mathrm{std}_m(\hat x_k^{(m)})$.

%======================================================================
\part{Distribution fitting}
%======================================================================

\section{Discretization and kernel}\label{sec:dist-kernel}

The spectrum is modeled as a continuous density $P(\BHF)$ (or $P(\dEQ)$)
discretized into $n$ \emph{bins} with equally spaced centers
$c_k\in[p_{\min},p_{\max}]$. The \emph{kernel}
\begin{equation}
  K_{ik} \;=\; A_{\text{sext}}^{(d=1)}\bigl(v_i;\;\BHF=c_k\bigr)
\end{equation}
is the absorption at $v_i$ of a sextet of unit depth with field $c_k$
(variable $\BHF$) or with $\dEQ=c_k$ and fixed $\BHF$ (variable $\dEQ$). The
kernel intensities follow $3{:}2{:}1$ by construction.

\section{Augmented linear system}\label{sec:dist-system}

The forward is
\begin{equation}
  \hat y_i \;=\; b + s\,v_i \;-\; \sum_{k}K_{ik}\,P_k \;-\;\sum_{m}K^{\text{sharp}}_{im}q_m,
\end{equation}
with $P_k\ge0$ the weights of the distribution and $q_m\ge0$ amplitudes of optional
``sharp'' sextets. Defining the design matrix
\begin{equation}
  X = \bigl[\;\bm 1\;\big|\;\bm v\;\big|\;{-}K\;\big|\;{-}K^{\text{sharp}}\;\bigr],
  \qquad \bm z=(b,s,\bm P^\top,\bm q^\top)^\top,
\end{equation}
and the weights $W=\diag(1/\sigma_i^2)$, the problem is
\begin{equation}
  \boxed{\;
  \min_{\bm z}\;\;\bigl\|W^{1/2}(X\bm z-\bm y)\bigr\|_2^2 \;+\; \alpha\,\mathcal R(\bm P)
  \quad\text{s.t.}\quad b\ge0,\;\bm P\ge0,\;\bm q\ge0.\;}
  \label{eq:dist-min}
\end{equation}
The slope $s$ has no bound. It is solved with bounded least squares
(\code{lsq\_linear}, reflective Newton with LSMR).

\subsection{Column-scale preconditioning}\label{sec:cond}
For small $\BHF$ the sextet collapses and the columns of $K$ become nearly
collinear with very disparate norms, which ill-conditions \eqref{eq:dist-min}.
Each column of the distribution is rescaled $K_k\to K_k/s_k$ with
$s_k=\lVert W^{1/2}K_k\rVert$ and $\tilde P_k=s_k P_k$; the penalizer and the weights
are adjusted accordingly and the scale is undone at the end. It is an
\emph{exact} reparametrization (solution and degrees of freedom invariant) that
only improves the numerical stability.

\section{Regularization}\label{sec:reg}

\subsection{Tikhonov (smoothness, L2)}
$\mathcal R(\bm P)=\lVert L\,\bm P\rVert_2^2$ with $L$ the \emph{second-difference}
operator $(n-2)\times n$:
\begin{equation}
  (L\bm P)_k = P_{k-1}-2P_k+P_{k+1},
\end{equation}
which penalizes the curvature and produces smooth distributions. The augmented
system stacks $\bigl[\begin{smallmatrix}W^{1/2}X\\ \sqrt\alpha\,L_{\text{ext}}\end{smallmatrix}\bigr]\bm z = \bigl[\begin{smallmatrix}W^{1/2}\bm y\\ \bm 0\end{smallmatrix}\bigr]$.

\subsection{Total variation (edges, L1)}
$\mathcal R(\bm P)=\lVert D\,\bm P\rVert_1$ with $D$ the \emph{first difference}
$(D\bm P)_k=P_{k+1}-P_k$. It favors piecewise-constant solutions that preserve
edges (less smearing of peaks). It is solved by reweighted least squares
(IRLS): starting from the majorization
$|t|\le t^2/(2|t^{\text{prev}}|)+\text{const}$, each iteration solves a weighted
Tikhonov with
\begin{equation}
  w_k = \frac{1}{\sqrt{(D\bm P^{\text{prev}})_k^2+\varepsilon^2}},
  \qquad \mathcal R\approx \sum_k w_k\,(D\bm P)_k^2,
\end{equation}
reusing \code{lsq\_linear}. The effective operator $\sqrt{w}\,D$ is also used
for the degrees of freedom and the error bars.

\section{Poisson likelihood in the distribution}
In Poisson mode the histogram fit is repeated with IRLS reweighting of
$\sigma$ from the model (\S\ref{sec:weighting}), closing the loop through
\code{lsq\_linear} until convergence.

\section{Effective degrees of freedom}\label{sec:effdof}

Regularization reduces the degrees of freedom below $n$. The effective
number is the trace of the hat matrix:
\begin{equation}
  p_{\text{eff}} \;=\; \tr A(\alpha),\qquad
  A(\alpha)=X\bigl(X^\top W X+\alpha L^\top L\bigr)^{-1}X^\top W,
  \label{eq:effdof}
\end{equation}
computable as $\tr\!\bigl[(X^\top WX+\alpha L^\top L)^{-1}X^\top WX\bigr]$. It is used
in $\chi^2_\nu=\chi^2/(N-p_{\text{eff}})$, avoiding the bias of counting the $n$
bins as independent parameters. It is invariant under the column rescaling
of \S\ref{sec:cond}.

\section{Selection of the regularization parameter}\label{sec:alpha}

Over a logarithmic grid of $\alpha$, three criteria are computed.

\paragraph{L-curve (maximum curvature).} With
$\rho(\alpha)=\log\lVert W^{1/2}(X\hat z-y)\rVert$ and
$\eta(\alpha)=\log\lVert L\hat{\bm P}\rVert$, the corner maximizes
\begin{equation}
  \kappa(\alpha)=\frac{\rho'\eta''-\eta'\rho''}{\bigl((\rho')^2+(\eta')^2\bigr)^{3/2}}.
\end{equation}

\paragraph{Trade-off.} Minimum normalized distance to the origin in the
$(\eta,\rho)$ plane.

\paragraph{Generalized cross-validation (GCV).}
\begin{equation}
  \mathrm{GCV}(\alpha)=\frac{\lVert W^{1/2}(\bm y - X\hat z)\rVert^2}{\bigl(N-p_{\text{eff}}(\alpha)\bigr)^2},
  \qquad \hat\alpha=\argmin_\alpha \mathrm{GCV}(\alpha),
\end{equation}
which does not require knowing the scale of the noise and is optimal in prediction loss.

\section{Error bars of \texorpdfstring{$P(\BHF)$}{P(BHF)}}\label{sec:perr}

For the regularized estimator $\hat{\bm z}=H^{-1}X^\top W\bm y$ with
$H=X^\top WX+\alpha L^\top L$, the linearized covariance is
\begin{equation}
  \widehat{\mathrm{Cov}}(\hat{\bm z}) \;=\; H^{-1}\,(X^\top WX)\,H^{-1}\cdot\frac{\chi^2}{N-p_{\text{eff}}},
\end{equation}
scaled by the effective $\chi^2_\nu$ (same convention as the discrete fit). The
square root of the diagonal corresponding to the bins gives $\sigma_{P_k}$. The band
$P_k\pm\sigma_{P_k}$ is clipped to $P\ge0$ when plotted (the non-negativity makes
it asymmetric in bins pinned to zero). In TV mode the effective operator
$\sqrt{w}\,D$ is used instead of $L$.

\section{Parametric distributions}
As an alternative to the regularized histogram, closed forms with few
parameters are fitted (nonlinear least squares):
\begin{align}
  \text{Gaussian:}\quad & P_k\propto \exp\!\bigl(-\tfrac12((c_k-\mu)/\varsigma)^2\bigr), & (b,s,A,\mu,\log\varsigma),\\
  \text{Binomial:}\quad & P_k\propto \binom{n-1}{k-1}p^{k-1}(1-p)^{n-k}, & (b,s,A,\mathrm{logit}\,p),\\
  \text{Fixed:}\quad & P_k=\phi_k\ \text{(preloaded)}, & (b,s,A).
\end{align}

\section{Global refinement (VARPRO)}\label{sec:refine}
When enabled, a nonlinear outer loop optimizes the shape parameters
$\bm\eta=(\delta_g,\log\gamma_g,\dots)$; at each evaluation the inner linear
problem \eqref{eq:dist-min} is solved \emph{entirely}, returning the
residual to the optimizer (separable VARPRO-type scheme):
\begin{equation}
  \min_{\bm\eta}\;\bigl\|W^{1/2}\bigl[X(\bm\eta)\,\hat{\bm z}(\bm\eta)-\bm y\bigr]\bigr\|^2,
  \qquad \hat{\bm z}(\bm\eta)=\argmin_{\bm z}\eqref{eq:dist-min}.
\end{equation}

%======================================================================
\appendix
\part{Appendices}

\section{Sextet calibration}\label{ap:calib}
The reference positions at $B_0=33\,$T are obtained from
$r^{(0)}=\tfrac12[-10.657,-6.167,-1.677,1.677,6.167,10.657]\cdot(B_0/32.95)\,$mm/s.
From them the Zeeman splittings used in \eqref{eq:Hkundig} are derived:
\begin{equation}
  \omega_e^{(0)} = r^{(0)}_2-r^{(0)}_1 \approx 2.248\ \text{mm/s},\qquad
  \omega_g^{(0)} = -\bigl[(r^{(0)}_4-r^{(0)}_3)+\omega_e^{(0)}\bigr] \approx -3.928\ \text{mm/s}.
\end{equation}

\section{Parameter bounds}\label{ap:bounds}
\begin{center}
\begin{tabular}{@{}lll@{}}
\toprule
Parameter & Bound & Unit\\
\midrule
$b$ (baseline) & $[0.7,\ 1.3]$ & --\\
$s$ (slope) & $[-5\cdot10^{-3},\ 5\cdot10^{-3}]$ & mm$^{-1}$s\\
$\delta$ & $[-2,\ 3]$ & mm/s\\
$\dEQ$ & $[0,\ 4]$ & mm/s\\
$\BHF$ & $[0,\ 50]$ & T\\
$\gamma_1$ & $[0.03,\ 2]$ & mm/s\\
$\gamma_2,\gamma_3$ & $[0.2,\ 3]$ & (relative)\\
$d$ (depth) & $[0,\ 0.30]$ & --\\
$i_1,i_2,i_3$ & $[0,9]/[0,6]/[0,3]$ & (relative)\\
$t$ (texture) & $[0,\ 1]$ & --\\
$\beta$ & $[0,\ 90]$ & degrees\\
\bottomrule
\end{tabular}
\end{center}

\section{Faddeeva function}
$w(z)=e^{-z^2}\erfc(-iz)$, evaluated by \code{scipy.special.wofz}. The real part
gives the Voigt profile \eqref{eq:voigt}; for $z=iy$ with $y\in\mathbb R$,
$w(iy)=e^{y^2}\erfc(y)$ is real, which allows the analytic normalization to the peak.

\section{Summary of the optimization stack}
\begin{center}
\begin{tabular}{@{}lll@{}}
\toprule
\textbf{Mode} & \textbf{Functional} & \textbf{Solver}\\
\midrule
Discrete & $\tfrac12\sum_i\rho(r_i^2)$ & \code{least\_squares} (TRF) [+ optional DE]\\
Tikhonov distribution & $\|W^{1/2}(X\bm z-\bm y)\|^2+\alpha\|L\bm P\|^2$ & \code{lsq\_linear} (bounded)\\
TV distribution & $\|W^{1/2}(X\bm z-\bm y)\|^2+\alpha\|D\bm P\|_1$ & IRLS over \code{lsq\_linear}\\
Parametric distribution & $\|W^{1/2}(X(\bm\theta)\bm z-\bm y)\|^2$ & \code{least\_squares} (NL)\\
Global refinement & $\|W^{1/2}(X(\bm\eta)\hat{\bm z}(\bm\eta)-\bm y)\|^2$ & outer NL + inner bounded\\
\bottomrule
\end{tabular}
\end{center}

\end{document}
